% =========================================================================
% Capítulo 3: Cinética de sistemas de partículas
% Tema 01: Sistemas de partículas
% Problema 01: Péndulo invertido sobre bloque móvil
% Origen: Ayudantía 2 (10), 2026
% =========================================================================
\subsubsection{Péndulo invertido sobre bloque móvil}

Un bloque de masa $M$ se desplaza sin fricción a lo largo de una superficie horizontal bajo la acción de una fuerza externa horizontal $F$. Una barra rígida sin masa de longitud $l$ está conectada al bloque mediante una bisagra en el punto $O$, permitiéndole rotar libremente en el plano vertical. En el extremo opuesto de la barra se encuentra una partícula de masa $m$, como se ilustra en la \textbf{Figura \ref{fig:pendulo_bloque}}. Sabiendo que la gravedad actúa hacia abajo y que el sistema parte de una configuración definida por el desplazamiento $x$ del bloque y el ángulo $\theta$ de la barra respecto a la vertical, realice lo siguiente:

\begin{enumerate}
    \item Defina los marcos de referencia y sistemas de coordenadas adecuados.
    \item Analice la cinemática de la partícula y del bloque para obtener sus vectores de posición, velocidad y aceleración.
    \item Analice la cinética del sistema para determinar el sistema de dos ecuaciones diferenciales de movimiento que describen el comportamiento de $x(t)$ y $\theta(t)$.
\end{enumerate}

\begin{figure}[!ht]
    \centering
    \begin{tikzpicture}[>=Latex, scale=1.3]
        % --- Parámetros ---
        \def\ang{35}      % Ángulo theta
        \def\Lbar{2.5}    % Longitud de la barra
        \def\Wblock{1.8}  % Ancho del bloque
        \def\Hblock{0.8}  % Alto del bloque

        % --- Suelo ---
        \fill[pattern=north east lines] (-2.5, -0.2) rectangle (4, 0);
        \draw[thick] (-2.5, 0) -- (4, 0);

        % --- Bloque M ---
        \draw[thick, fill=gray!40] (0, 0.2) rectangle (\Wblock, \Hblock+0.2);
        \node at (\Wblock/2, \Hblock/2+0.2) {$M$};
        
        % Ruedas
        \draw[fill=white] (0.4, 0.1) circle (0.1);
        \draw[fill=white] (\Wblock-0.4, 0.1) circle (0.1);

        % Punto O (Bisagra)
        \coordinate (O) at (\Wblock/2, \Hblock+0.2);
        \filldraw[fill=gray!20] (O) circle (2pt) node[above left, xshift=-2pt] {$O$};

        % --- Barra y Partícula m ---
        \begin{scope}[shift={(O)}, rotate=-\ang]
            \draw[line width=4pt, line cap=round] (0,0) -- (0, \Lbar) node[midway, right] {$l$};
            \draw[line width=3pt, gray, line cap=round] (0,0) -- (0, \Lbar);
            \shadedraw[ball color=gray!30] (0, \Lbar) circle (0.2) node[above right=2pt] {$P$} node[font=\footnotesize] {$m$};
            \coordinate (P) at (0, \Lbar);
        \end{scope}
        \filldraw[] (O) circle (0.6pt);

        % --- Fuerzas y Referencias ---
        % Fuerza F
        \draw[->, ultra thick] (-1, \Hblock/2+0.2) -- (0, \Hblock/2+0.2) node[midway, above left, xshift=-5pt] {$\mathbf{F}$};

        % Eje x de referencia
        \draw[|->, thick] (\Wblock/2, \Hblock+2) -- (\Wblock/2 + 1, \Hblock+2) node[midway, above] {$x$};
        
        % Ángulo theta
        \draw[dashed, gray] (O) -- (\Wblock/2, \Hblock+1.5);
        \draw[->] (\Wblock/2, \Hblock+1.0) arc (90:{90-\ang}:0.8);
        \node at (\Wblock/2+0.3, \Hblock+1.3) {$\theta$};

        % Vector Gravedad
        \draw[->, thick] (-1.5, 2.5) -- (-1.5, 1.5) node[midway, left] {$\mathbf{g}$};

    \end{tikzpicture}
    \caption{Diagrama de una partícula de masa $m$ acoplada a un bloque de masa $M$ con movimiento horizontal y rotación libre.}
    \label{fig:pendulo_bloque}
\end{figure}

% --- SOLUCIÓN ---------------------------------------------------------------
\begin{solucion}{Péndulo invertido sobre bloque móvil}
\textbf{Marcos de referencia y Sistemas de Coordenadas}
\begin{itemize}
    \item $\mathcal{F}$: inercial, fijo al suelo, con origen en la posición inicial de la bisagra $O$ $\longrightarrow\{\vec{E}_x;\vec{E}_y;\vec{E}_z\}$, con $\vec{E}_x$ horizontal hacia la derecha, $\vec{E}_y$ vertical hacia arriba y $\vec{E}_z=\vec{E}_x\times\vec{E}_y$ saliendo de la página.
    \item $\mathcal{A}$: solidario a la varilla $OP$ $\longrightarrow\{\vec{e}_r;\vec{e}_\theta;\vec{e}_z\}$, con $\vec{e}_r$ a lo largo de $OP$, $\vec{e}_z=\vec{E}_z$ y $\vec{e}_\theta=\vec{e}_z\times\vec{e}_r$.
\end{itemize}

\begin{center}
\begin{tikzpicture}[>=Latex, scale=1.3]
    % --- PARÁMETROS ---
    \def\ang{35}   % ángulo theta [°]
    \def\len{1.5}  % largo de los vectores [u]

    % --- BASE FIJA ---
    \draw[->, thick, azul] (0,0) -- (\len,0) node[right] {$\vec{E}_x$};
    \draw[->, thick, azul] (0,0) -- (0,\len) node[above] {$\vec{E}_y$};
    \draw[dashed, gray] (0,0) -- (0,-0.3);

    % --- BASE MÓVIL ---
    \draw[->, thick, myred] (0,0) -- ({90-\ang}:\len) node[above right] {$\vec{e}_r$};
    \draw[->, thick, myred] (0,0) -- ({180-\ang}:\len) node[above left] {$\vec{e}_\theta$};

    % --- ÁNGULO ---
    \draw[->] (0,0.8) arc (90:{90-\ang}:0.8);
    \node at ({90-\ang/2}:1.05) {$\theta$};
    \filldraw (0,0) circle (1.2pt) node[below left] {$O$};
\end{tikzpicture}
\end{center}

De la geometría de la Figura (donde $\theta$ se mide en sentido horario desde la vertical superior $\vec{E}_y$), la relación entre las bases y la velocidad angular de la varilla son:
\begin{equation*}
\begin{split}
    \vec{e}_r &= \sin\theta\,\vec{E}_x + \cos\theta\,\vec{E}_y\\
    \vec{e}_\theta &= -\cos\theta\,\vec{E}_x + \sin\theta\,\vec{E}_y
\end{split}
\end{equation*}
$$\velang{F}{A} = -\dot{\theta}\,\vec{E}_z$$

\subsubsection*{Cinemática del sistema de partículas}
\textbf{Bloque de masa $M$ (punto $O$)}
$$\vec{r}_M = x\,\vec{E}_x \quad\Longrightarrow\quad \marc{\vec{v}_M}{F} = \dot{x}\,\vec{E}_x \quad\Longrightarrow\quad \marc{\vec{a}_M}{F} = \ddot{x}\,\vec{E}_x$$

\textbf{Partícula de masa $m$ (punto $P$)}\\
La posición relativa es $\vec{r}_{P/O} = l\,\vec{e}_r$. Aplicando el Teorema de Transporte dos veces:
$$\marc{\vec{v}_m}{F} = \marc{\vec{v}_M}{F} + \velang{F}{A}\times\vec{r}_{P/O} = \dot{x}\,\vec{E}_x + (-\dot{\theta}\,\vec{E}_z) \times (l\,\vec{e}_r) = \dot{x}\,\vec{E}_x - l\dot{\theta}\,\vec{e}_\theta$$
$$\marc{\vec{a}_m}{F} = \marc{\vec{a}_M}{F} + \der{(-l\dot{\theta}\,\vec{e}_\theta)}{A}{} + \velang{F}{A} \times (-l\dot{\theta}\,\vec{e}_\theta) = \ddot{x}\,\vec{E}_x - l\ddot{\theta}\,\vec{e}_\theta - l\dot{\theta}^2\,\vec{e}_r$$
Sustituyendo $\vec{e}_r$ y $\vec{e}_\theta$ en la base inercial $\{\vec{E}_x, \vec{E}_y\}$:
$$\marc{\vec{a}_m}{F} = \left(\ddot{x} + l\ddot{\theta}\cos\theta - l\dot{\theta}^2\sin\theta\right)\vec{E}_x + \left(-l\ddot{\theta}\sin\theta - l\dot{\theta}^2\cos\theta\right)\vec{E}_y$$

\subsubsection*{Cinética del sistema de partículas}
\textbf{Primera ecuación: Newton para el conjunto $\{M, m\}$ en $\vec{E}_x$}\\
Para el sistema completo $\{M, m\}$, la fuerza axial interna de la varilla se anula por acción y reacción. La única fuerza externa con componente horizontal es $\vec{F} = F\,\vec{E}_x$. Aplicando $\sum \vec{F}_{\text{ext}} \cdot \vec{E}_x = M a_{M,x} + m a_{m,x}$:
$$F = M\ddot{x} + m\left(\ddot{x} + l\ddot{\theta}\cos\theta - l\dot{\theta}^2\sin\theta\right)$$
\respuesta{$\displaystyle (M + m)\ddot{x} + ml\ddot{\theta}\cos\theta - ml\dot{\theta}^2\sin\theta = F$}

\textbf{Segunda ecuación: Newton para $m$ en dirección transversal $\vec{e}_\theta$}\\
Aislando la partícula $m$, las fuerzas que actúan sobre ella son la reacción axial de la varilla $\vec{R} = R\,\vec{e}_r$ y el peso $-mg\,\vec{E}_y$. Proyectando sobre $\vec{e}_\theta$ eliminamos la fuerza de reacción desconocida $R$:
$$\sum \vec{F}_m \cdot \vec{e}_\theta = m\left(\marc{\vec{a}_m}{F}\cdot \vec{e}_\theta\right)$$
Dado que $(-mg\,\vec{E}_y) \cdot \vec{e}_\theta = -mg\sin\theta$ y $\vec{E}_x \cdot \vec{e}_\theta = -\cos\theta$:
$$-mg\sin\theta = m\left(\ddot{x}\,\vec{E}_x - l\ddot{\theta}\,\vec{e}_\theta - l\dot{\theta}^2\,\vec{e}_r\right) \cdot \vec{e}_\theta = m\left(-\ddot{x}\cos\theta - l\ddot{\theta}\right)$$
Dividiendo por $-m$ y ordenando términos, se obtiene la segunda ecuación diferencial:
\respuesta{$\displaystyle l\ddot{\theta} + \ddot{x}\cos\theta - g\sin\theta = 0$}

\patoojo{Proyecta sobre la dirección perpendicular a la fuerza que no conoces. Aquí la varilla solo empuja en $\vec{e}_r$, así que al proyectar la ecuación de $m$ sobre $\vec{e}_\theta$ desaparece la reacción $R$ y no hay que despejarla.}

\patotip{Fíjate en el signo de $g\sin\theta$: es un péndulo \emph{invertido} ($\theta$ medido desde la vertical hacia arriba), por eso la gravedad desestabiliza y el término aparece como $-g\sin\theta$, con el signo contrario al $+g\sin\theta$ del péndulo común.}
\end{solucion}
